% =====================================================================
%  FICHE 4 — THEME 1 — Corrige niveau FACILE
% =====================================================================
\documentclass[11pt,a4paper]{article}
\usepackage[utf8]{inputenc}
\usepackage[T1]{fontenc}
\usepackage{lmodern}
\usepackage[french]{babel}
\usepackage{amsmath,amssymb,mathtools}
\usepackage{icomma}
\usepackage[version=4]{mhchem}
\usepackage{siunitx}
\sisetup{output-decimal-marker={,},per-mode=symbol}
\DeclareSIUnit\Molar{M}
\usepackage{xcolor}
\usepackage{tikz}
\usepackage[most]{tcolorbox}
\usepackage{enumitem}
\usepackage{booktabs}
\usepackage{array}
\usepackage{fancyhdr}
\usepackage[a4paper,margin=2.2cm,headheight=26pt,headsep=10pt]{geometry}

\definecolor{primary}{RGB}{146,95,160}
\definecolor{primarydark}{RGB}{92,52,104}
\definecolor{excol}{RGB}{0,135,90}

\tcbset{boxbase/.style={enhanced,breakable,boxrule=0.9pt,arc=2.5pt,
   left=3mm,right=3mm,top=2mm,bottom=2mm,fonttitle=\bfseries,coltitle=white,
   toptitle=0.5mm,bottomtitle=0.5mm}}

\newcounter{exo}
\newenvironment{corr}[1][]{%
  \par\medskip\refstepcounter{exo}%
  \noindent\begin{tcolorbox}[boxbase,colback=excol!5,colframe=excol,
     title={\textsc{Correction \theexo}\ifx\relax#1\relax\relax\else\textmd{ --- \itshape #1}\fi}]}%
  {\end{tcolorbox}}

\pagestyle{fancy}\fancyhf{}
\fancyhead[L]{\small\textcolor{primarydark}{\textbf{Fiche 4} --- Th.\,1 : Mole, masse molaire, entit\'es}}
\fancyhead[R]{\small\textcolor{primarydark}{\textsc{Chimie} --- \textbf{Corrigé Facile}}}
\fancyfoot[C]{\small\thepage}
\renewcommand{\headrulewidth}{0.8pt}
\renewcommand{\headrule}{\hbox to\headwidth{\color{primary}\leaders\hrule height \headrulewidth\hfill}}
\newcommand{\NA}{N_{\!\text{A}}}
\setlength{\parindent}{0pt}

\begin{document}

\begin{center}
{\LARGE\bfseries\color{primarydark} Thème 1 --- Corrigé Facile}
\end{center}

\begin{corr}[masse molaire d'un composé]
On additionne les masses molaires atomiques, pondérées par le nombre d'atomes de la formule.
\[ \textbf{a)}\quad M(\ce{CO2})=12+2(16)=\SI{44}{\gram\per\mole}. \]
\[ \textbf{b)}\quad M(\ce{C6H12O6})=6(12)+12(1)+6(16)=72+12+96=\SI{180}{\gram\per\mole}. \]
\[ \textbf{c)}\quad M(\ce{CH4N2O})=12+4(1)+2(14)+16=12+4+28+16=\SI{60}{\gram\per\mole}. \]
\[ \textbf{d)}\quad M(\ce{FeSO4})=56+32+4(16)=56+32+64=\SI{152}{\gram\per\mole}. \]
\textit{Réflexe :} recompter chaque atome (le glucose a bien $6+6=12$ atomes d'\ce{O}\ldots non :
$6$ atomes d'\ce{O} ; les $12$ concernent \ce{H}).
\end{corr}

\begin{corr}[de la masse à la quantité de matière]
On part d'une masse et l'on cherche des moles : on \textbf{divise} par la masse molaire.
\[ n=\frac{m}{M}=\frac{20}{40}=\SI{0.50}{\mole}. \]
\[ \boxed{n(\ce{NaOH})=\SI{0.50}{\mole}} \]
\textit{Vérification des unités :} $\dfrac{\si{\gram}}{\si{\gram\per\mole}}=\si{\mole}$. Cohérent.
\end{corr}

\begin{corr}[de la quantité de matière au nombre d'entités]
\textbf{a)} On part des moles et l'on veut des entités : on \textbf{multiplie} par $\NA$.
\[ N=n\times\NA=0{,}25\times\num{6.02e23}=\num{1.505e23}\ \text{molécules}. \]
\[ \boxed{N\approx\num{1.51e23}\ \text{molécules d'eau}} \]
\textbf{b)} Chaque molécule \ce{H2O} contient \textbf{2} atomes d'hydrogène ; la quantité d'atomes
d'\ce{H} est donc le double de celle des molécules :
\[ n(\ce{H})=2\times n(\ce{H2O})=2\times 0{,}25=\SI{0.50}{\mole}\ \text{d'atomes d'hydrogène}. \]
\textit{Piège classique :} confondre « moles de molécules » et « moles d'atomes ».
\end{corr}

\begin{corr}[du nombre d'entités à la quantité de matière]
\textbf{a)} On part d'un nombre d'entités et l'on veut des moles : on \textbf{divise} par $\NA$.
\[ n=\frac{N}{\NA}=\frac{\num{3.01e23}}{\num{6.02e23}}=\SI{0.500}{\mole}. \]
\[ \boxed{n(\ce{CO2})=\SI{0.500}{\mole}} \]
(C'est cohérent : $\num{3.01e23}$ est la moitié de $\NA$, donc une demi-mole.)\\[2pt]
\textbf{b)} Chaque molécule \ce{CO2} contient \textbf{2} atomes d'oxygène :
\[ n(\ce{O})=2\times 0{,}500=\SI{1.00}{\mole}\ \text{d'atomes d'oxygène}. \]
\end{corr}

\begin{corr}[masse molaire d'un hydrate]
On calcule la masse molaire de la partie « sel » puis on ajoute $x\times\SI{18}{\gram\per\mole}$
pour les $x$ molécules d'eau.
\[ \textbf{a)}\quad M(\ce{CuSO4.5H2O})=\underbrace{63{,}5+32+4(16)}_{M(\ce{CuSO4})=159{,}5}+5\times 18
   =159{,}5+90=\SI{249.5}{\gram\per\mole}. \]
\[ \textbf{b)}\quad M(\ce{Na2CO3.10H2O})=\underbrace{2(23)+12+3(16)}_{M(\ce{Na2CO3})=106}+10\times 18
   =106+180=\SI{286}{\gram\per\mole}. \]
\textit{À retenir :} l'eau de cristallisation compte pleinement dans la masse molaire de l'hydrate.
\end{corr}

\end{document}
